\subsection{Alínea 3.t) --- Functor $F: \mathbf{Turing\ Machines} \to \mathbf{Struct}$}
\label{subsec:functor-turing-struct}
\subsubsection{Functor Covariante $F: \mathbf{Turing\ Machines} \to \mathbf{Struct}$}

\begin{definition}[Functor covariante $F: \mathbf{Turing\ Machines} \to \mathbf{Struct}$]
Dada uma MT $M=(Q,\Gamma,\delta,q_0,F_{\mathrm{acc}},B)$ com $|Q|=m$, define-se o triplo $(d,\varphi,m)$:
\begin{itemize}
    \item $\varphi(q) = \delta(q, B)_1$ — estado de transição ao ler o símbolo em branco, capturando a dinâmica sob um símbolo fixo;
    \item $d(q) = \mathbf{1}_{q \in F_{\mathrm{acc}}} + \mathbf{i}\cdot\mathbf{1}_{q=q_0}$, distinguindo os quatro tipos de estados em $\{0,\,1,\,\mathbf{i},\,1{+}\mathbf{i}\}$.
\end{itemize}
Dado $(f_Q, f_\Gamma): M_1 \to M_2$, define-se $F(f_Q,f_\Gamma) = f_Q$.
\end{definition}

\begin{proposition}
$F$ é covariante e bem definido: $f_\Gamma(B_1)=B_2$ implica $f_Q \circ \varphi_1 = \varphi_2 \circ f_Q$; a preservação de $q_0$ e $F_{\mathrm{acc}}$ garante $d_1 = d_2 \circ f_Q$.
\end{proposition}

\begin{lstlisting}[style=matlab, caption={Functor covariante $F: \mathbf{Turing\ Machines} \to \mathbf{Struct}$}]
function S = functor_TM_to_Struct(M)
    S.n   = M.Q;
    S.phi = M.delta(:, M.blank, 1)';  % phi(q) = delta(q, blank)_1
    S.d   = zeros(1, M.Q, 'like', 1i);
    for q = 1:M.Q
        S.d(q) = double(ismember(q,M.F)) + 1i*double(q==M.q0);
    end
end

function fS = functor_TM_to_Struct_morphism(fQ)
    fS = fQ;
end
\end{lstlisting}

\subsubsection{Functor Contravariante --- Análise de Existência}

\begin{proposition}[Inexistência de functor contravariante não trivial]
Não existe functor contravariante canónico $G: \mathbf{Turing\ Machines} \to \mathbf{Struct}$. Um morfismo $(f_Q,f_\Gamma): M_1 \to M_2$ induziria $G(f_Q,f_\Gamma): G(M_2) \to G(M_1)$, ou seja, um morfismo de \textbf{Struct} com $g \circ \varphi_2 = \varphi_1 \circ g$ e $d_2 = d_1 \circ g$. A condição sobre $d$ requereria $g(q_0^{(2)}) = q_0^{(1)}$ e $g(F_2) \subseteq F_1$, que é exatamente uma condição de morfismo ``inverso'' em \textbf{TM}. Tal mapa $g$ só existe canonicamente quando $f_Q$ é invertível.
\end{proposition}

\begin{remark}
Na subcategoria dos isomorfismos de \textbf{Turing\ Machines}, $G(f_Q,f_\Gamma) = F(f_Q^{-1}, f_\Gamma^{-1})$ define um functor contravariante. No caso geral, só existe o functor contravariante constante $G(M) = S_0$, $G(f) = \mathrm{id}_{S_0}$.
\end{remark}

\subsubsection{Transformação Natural}

Seja $G: \mathbf{Turing\ Machines} \to \mathbf{Struct}$ um segundo functor que usa um símbolo de referência diferente $a^{**} \neq a^*$ para definir o endomapa $\psi(q) = \delta(q, a^{**})_1$, mas o mesmo mapa $d$. Então existe uma transformação natural $\eta: F \Rightarrow G$ cujas componentes são:
\[
    \eta_M : F(M) \to G(M), \qquad \eta_M = \mathrm{id}_{\{1,\ldots,m\}},
\]
desde que $\varphi$ e $\psi$ coincidam (o que acontece se $\delta(q, a^*) = \delta(q, a^{**})$ para todo $q$). No caso geral, $\eta_M$ é a identidade mas deixa de ser morfismo de \textbf{Struct} (pois $\varphi \neq \psi$), pelo que a transformação natural genuína requer que os dois functores partilhem o mesmo símbolo de referência.

Uma transformação natural sempre definida é $\varepsilon: F \Rightarrow C_{S_0}$, onde $C_{S_0}$ é o functor constante no objeto trivial $S_0 = (0, \mathrm{id}, 1)$, com $\varepsilon_M$ o único morfismo de \textbf{Struct} para $S_0$.

\begin{lstlisting}[style=matlab, caption={Transformação natural $\eta: F \Rightarrow G$ (mudança de símbolo de referência)}]
function ok = check_natural_transformation_TM_Struct(M, sym1, sym2)
    % Verifica se id_Q e morfismo de Struct entre F(M) usando sym1
    % e G(M) usando sym2, i.e., se phi1 == phi2
    phi1 = zeros(1, M.Q);
    phi2 = zeros(1, M.Q);
    for q = 1:M.Q
        phi1(q) = squeeze(M.delta(q, sym1, :))'  ;
        phi1(q) = M.delta(q, sym1, 1);
        phi2(q) = M.delta(q, sym2, 1);
    end
    ok = isequal(phi1, phi2);
    % Se ok=true, eta_M = identidade e natural_transformacao existe
end
\end{lstlisting}